A linear probe can reveal a clinical label without showing that its decoded direction selectively affects a model's answer. We test this distinction at the exact final visual block consumed downstream by LLaVA-1.5-7B and Qwen2.5-VL-7B. Our registered protocol combines a fixed 512-dimensional readout, patient-disjoint evaluation, type-to-random-label controls, patient-bootstrap uncertainty, and relative-token probe-normal interventions against random, sham, and fixed clinical directions. Across three model--concept cells on NIH ChestX-ray14, AUROC ranges from 0.7742 to 0.8009 and every selectivity interval excludes zero. Neither LLaVA Effusion nor LLaVA Edema clears random/sham controls. Qwen Effusion clears them twice, changing mean $P(yes)$ by 0.1945 on 400 independent patients, but Nodule changes the same answer by 0.2519; the familywise Effusion-minus-maximum-unrelated margin is $-0.0573$ (95\% CI $[-0.0694,-0.0450]$). Two prospective mechanism gates sharpen this boundary. A six-direction by six-question Qwen matrix yields zero owned concepts and no shared alias beyond 119 random directions. A fresh 50-patient Consolidation test finds positive controlled selectivity but no reliable input-displacement closure: the concept gain is 0.0022 versus a random maximum of 0.0025. Thus, in these registered cells, linear availability at a consumed visual locus does not establish concept-specific causal use under probe-normal intervention. The experiments separate decodability, generic perturbation sensitivity, clinical direction specificity, and mechanism confirmation.
